\section{Pre-registration record}\label{app:prereg}

\paragraph{Rule.} Each analysis was registered before it was run. Before any estimate, its
protocol was committed and pushed. The protocol fixes:
\begin{itemize}
  \item the data and sample;
  \item the estimator and its settings;
  \item the decision rules and multiple-comparison families;
  \item the sensitivity analyses;
  \item for the case study, the hypotheses that cannot be tested reliably.
\end{itemize}
Table~\ref{tab:prereg} lists the commit that first added each protocol. In every case the
commit precedes, in the repository history, the commit that first added the results.

\paragraph{Changes after registration.} A change made after registration is recorded in
the protocol's change log, with its timing relative to the results and its reason. Three
kinds of change occurred:
\begin{itemize}
  \item \emph{clarifications made before any estimate}: for example, the Bayesian claim
    rules when a specification has no population mean, and the treatment of an indicator
    without variation;
  \item \emph{corrections after the first computation that change no rule}: for example,
    the coverage of employment in persons;
  \item \emph{amendments}: registered extensions with their own decision rules, written
    before their fits. Examples are the dynamic and regime Bayesian models and the monthly
    identification.
\end{itemize}
Analyses added after results were seen, such as the checks of the common factor, are
labelled as exploratory in the protocol and in the report.

\paragraph{The classical plan.} \texttt{METHODS.md} is the research plan of the classical
stage. It was committed with the data pipeline before any estimate and revised in later
commits without a separate change log; the repository history records each revision.

% Generated by eu-okun publication. Do not edit.
\begin{table}[t]\centering\small
\caption{Pre-registration record: the commit that first added each protocol (every one precedes the commit that added its results), why the analysis exists, and the number of top-level entries in its change log. Initial plan: an item of the first research roadmap, committed before any estimate. Follow-up: written after the results of an earlier stage (results) or in response to an external review (review), as the protocol states; its rules were fixed before its own estimates, but its question was chosen knowing earlier results.}
\label{tab:prereg}
\footnotesize\setlength{\tabcolsep}{4pt}
\begin{tabular}{lllr}
\toprule
Protocol & First commit & Origin & Change log \\
\midrule
\texttt{METHODS.md} & \texttt{00f0ada}, 2026-10-14 & initial plan & 0 \\
\texttt{SMALL\_\allowbreak{}SAMPLE\_\allowbreak{}PROTOCOL.md} & \texttt{6c8c998}, 2026-10-17 & follow-up: results & 2 \\
\texttt{HEADLINE\_\allowbreak{}INFERENCE\_\allowbreak{}PROTOCOL.md} & \texttt{6f0d8bd}, 2026-10-19 & follow-up: review & 1 \\
\texttt{CORE\_\allowbreak{}GRID\_\allowbreak{}PROTOCOL.md} & \texttt{84f50e1}, 2026-10-19 & follow-up: review & 0 \\
\texttt{SMOOTHING\_\allowbreak{}PROTOCOL.md} & \texttt{e0edaef}, 2026-10-19 & follow-up: review & 1 \\
\texttt{GDP\_\allowbreak{}MEASUREMENT\_\allowbreak{}PROTOCOL.md} & \texttt{e2d577f}, 2026-10-19 & follow-up: review & 0 \\
\texttt{HOURS\_\allowbreak{}PROTOCOL.md} & \texttt{55783e3}, 2026-10-20 & follow-up: review & 2 \\
\texttt{ASYMMETRY\_\allowbreak{}PROTOCOL.md} & \texttt{39b6b00}, 2026-10-21 & follow-up: review & 1 \\
\texttt{CCE\_\allowbreak{}PROTOCOL.md} & \texttt{d1f4c13}, 2026-10-21 & follow-up: review & 0 \\
\texttt{AGGREGATE\_\allowbreak{}PROTOCOL.md} & \texttt{474c22b}, 2026-10-21 & follow-up: review & 0 \\
\texttt{MEASUREMENT\_\allowbreak{}PROTOCOL.md} & \texttt{0d19409}, 2026-10-22 & follow-up: review & 0 \\
\texttt{BAYES\_\allowbreak{}PROTOCOL.md} & \texttt{24c4704}, 2026-10-15 & initial plan; amendments: results & 10 \\
\texttt{BAYES\_\allowbreak{}DIAGNOSTICS\_\allowbreak{}PROTOCOL.md} & \texttt{8265ab5}, 2026-10-22 & follow-up: review & 1 \\
\texttt{PREDICTIVE\_\allowbreak{}PROTOCOL.md} & \texttt{109eee5}, 2026-10-16 & initial plan & 1 \\
\texttt{BAYES\_\allowbreak{}FORECAST\_\allowbreak{}PROTOCOL.md} & \texttt{3f9a2b7}, 2026-10-16 & initial plan & 0 \\
\texttt{IDENTIFICATION\_\allowbreak{}PROTOCOL.md} & \texttt{7d01063}, 2026-10-16 & initial plan & 3 \\
\texttt{MONTHLY\_\allowbreak{}IDENTIFICATION\_\allowbreak{}PROTOCOL.md} & \texttt{97fc26b}, 2026-10-18 & follow-up: results & 3 \\
\texttt{SHIFT\_\allowbreak{}SHARE\_\allowbreak{}PROTOCOL.md} & \texttt{46ca0a2}, 2026-10-18 & follow-up: results & 1 \\
\texttt{SIGN\_\allowbreak{}RESTRICTIONS\_\allowbreak{}PROTOCOL.md} & \texttt{8410575}, 2026-10-18 & follow-up: results & 3 \\
\texttt{HETEROGENEITY\_\allowbreak{}PROTOCOL.md} & \texttt{21a4e23}, 2026-10-18 & initial plan & 2 \\
\bottomrule
\end{tabular}
\end{table}

